\documentclass[10pt,a4paper]{amsart}
\usepackage[margin=19mm]{geometry}
\usepackage{amsmath,amssymb,mathtools}
\usepackage[scr=rsfs,cal=euler]{mathalfa}
\usepackage{xfrac}
\usepackage[unicode,hidelinks,bookmarks=false]{hyperref}
\newcommand{\ie}{i.e.\ }
\newcommand{\cf}{cf.\ }
\newcommand{\resp}{respectively\ }
\newcommand{\Subsection}{\subsection}
\providecommand{\nobreakdash}{\nobreak\hbox{-}}
\setlength{\emergencystretch}{2em}
\multlinegap0pt
\usepackage{mathtools}
\usepackage{placeins}
\usepackage{cleveref}
\usepackage{enumitem}
\hypersetup{colorlinks=true,linkcolor=blue,citecolor=blue,urlcolor=blue}
\newtheorem{theorem}{Theorem}[section]
\newtheorem{proposition}[theorem]{Proposition}
\newtheorem{corollary}[theorem]{Corollary}
\newtheorem{lemma}[theorem]{Lemma}
\newtheorem{conjecture}[theorem]{Conjecture}
\newtheorem{definition}[theorem]{Definition}
\newtheorem{example}[theorem]{Example}
\newtheorem{problem}[theorem]{Problem}
\newtheorem{remark}[theorem]{Remark}
\newcommand{\C}{\mathbb C}
\newcommand{\R}{\mathbb R}
\newcommand{\CP}{\mathbb {CP}}
\newcommand{\RP}{\mathbb {RP}}
\newcommand{\I}{\mathfrak I}
\newcommand{\ord}{\operatorname{ord}}
\newcommand{\Res}{\operatorname{Res}}
\renewcommand{\Re}{\operatorname{Re}}
\renewcommand{\Im}{\operatorname{Im}}
\begin{document}
\title[Inflection curves of rational vector fields]{Inflection curves of rational vector fields}
\author{Boris Shapiro, Guillaume Tahar}
\date{}
\maketitle
\noindent\emph{Source and licence.} Inflection curves of rational vector fields. Version arXiv:2605.28266v1 (CC BY 4.0). This material is distributed under the Creative Commons Attribution 4.0 International licence, http://creativecommons.org/licenses/by/4.0/, as recorded in the arXiv OAI licence metadata for this exact version and shown on the abstract page https://arxiv.org/abs/2605.28266v1. Full rights notice: licenses/arxiv-2605.28266-CC-BY-4.0.txt.\par\smallskip
\noindent\textit{Editorial scope.} Counted unit: the original \texttt{theorem} environment \texttt{thm:local-pole} at source lines 199--211 together with its complete proof at lines 213--227; the delivered document keeps source lines 197--228 so that every referenced label and every citation resolves inside it. Native editable source is the paper's own concatenated TeX; the AMS wrapper is editorial formatting only. No publisher class, upstream script or unrelated artwork is redistributed.
\section{Local structure near poles and ordinary singularities}


\begin{theorem}\label{thm:local-pole}
Suppose that $z_0$ is a pole of $R$ of order $s\ge1$.  Write
\[
        R(z)=\frac{a}{(z-z_0)^s}+O((z-z_0)^{-s+1}),
        \qquad a\ne0.
\]
Then near $z_0$, the inflection curve $\I_R$ has $s+1$ smooth real analytic branches through $z_0$, with tangent rays determined by
\begin{equation}\label{eq:pole-directions}
        \arg(z-z_0)=\frac{\arg(-sa)+m\pi}{s+1},
        \qquad m=0,1,\ldots,2s+1.
\end{equation}
Opposite rays in \eqref{eq:pole-directions} belong to the same analytic branch.  In particular, if $z_0$ is a simple pole of $R$, then $z_0$ is an ordinary real node of $\I_R$.
\end{theorem}

\begin{proof}
Put $w=z-z_0$.  Then
\[
        R'(z)=-sa w^{-s-1}+O(w^{-s}).
\]
The equation $\Im R'(z)=0$ is equivalent, after multiplication by $|w|^{2s+2}$, to
\[
        \Im\bigl((-sa)\overline w^{s+1}+O(|w|^{s+2})\bigr)=0.
\]
The leading homogeneous part is $\Im((-sa)\overline w^{s+1})$.  Writing $w=re^{i\theta}$, its zero set is
\[
        \arg(-sa)-(s+1)\theta\in \pi\mathbb Z,
\]
which is equivalent to \eqref{eq:pole-directions}.  The leading homogeneous polynomial has $2s+2$ distinct real linear factors.  The standard real-analytic factorization of plane curve germs with distinct tangent cone therefore gives $s+1$ smooth analytic branches.
\end{proof}


\end{document}
